\scanpage{9}
% Pencil foliation: 8, upper right and lower left.
Alors on trouve un $\otimes$-foncteur \inserted{pleinement} fidèle
\diagram{\begin{array}{c}
\pi_1\text{-}\mathbb{Q}\text{-modules admissibles de rang fini sur }\mathbb{Q}\\
\longrightarrow\quad\text{motifs sur }S
\end{array}}
qui est décrit par un \inserted{épi}morphisme
\diagram{\begin{tikzcd}
G=\varprojlim G_\alpha \arrow[r] & \pi_1(S,\unsure{\xi|\zeta})
\end{tikzcd}}
et se factorise donc en un \uline{épimorphisme} canonique
\margin{$*$}
\diagram{\begin{array}{|ccc|}\hline
\varprojlim_\alpha G_\alpha/G_\alpha^0&\longrightarrow&\pi_1(S,\unsure{\xi|\zeta})\\\hline
\end{array}}
Soit $\pi'_1$ le premier mb. Alors \deleted{les}
\inserted{(supposant $K=\mathbb{Q}$} \deleted{$\pi'_1$}
pour simplifier) les $\pi'_1$-$\mathbb{Q}$-modules
correspondent à la sous-catégorie de $\mathcal{M}(S)$ formée
des motifs \inserted{$M$} tels que la sous-\inserted{$\otimes$-}catégorie engendrée
par $M$ soit de ``longueur finie''. Utilisant
conjecture de Tate (ou de Hodge si $S$ de
car. générique nulle) on \inserted{peut} voi\deleted{t}\inserted{\unsure{r}} que ce sont les
motifs de \deleted{\unsure{poids}} niveau \inserted{et poids} zéro, \inserted{i.e. les motifs effectifs de poids zéro,}
i.e. que \uline{c'est un
isomorphisme}. Les \deleted{résultats}\inserted{réflexions} précédentes peuvent
se \deleted{\unsure{faire}}\inserted{développer} aussi (utilisant la notion de
motif sur $S$ à coeff. dans $K$) si on ne suppose
pas $K=\mathbb{Q}$.

\deleted{\unsure{NB} \illegible{2} \unsure{objets} de \unsure{$\alpha$} \unsure{que}}

\uline{NB} Soit $S'\longrightarrow S$ un morphisme
étale \inserted{\unsure{surj.}}\deleted{de} $S$, \inserted{d'avance,} et supposons que $F$ provienne d'

% Printed letterhead, upside down relative to the handwriting, lower right.
INSTITUT

DES HAUTES ETUDES SCIENTIFIQUES

91 - BURES-SUR-YVETTE

928 48-53
